\subsection{张量积与同态模的分式模}

命题 18. 设 A 是环，S 是 A 的乘性子集。

\begin{enumerate}
\def\labelenumi{(\roman{enumi})}
\tightlist
\item
  若 M 和 N 是两个 A-模，则 \(S^{-1}A\)-模 ( \(S^{-1}M\) ) \(\otimes_{A}N\)、\(M\otimes_{A}(S^{-1}N)\)、( \(S^{-1}M\) ) \(\otimes_{S^{-1}A}(S^{-1}N)\) 与 \(S^{-1}(M\otimes_{A}N)\) 典范地同构。
\item
  若 \(M'\) 与 \(N'\) 是两个 \(S^{-1}A\)-模，则典范同态

\[
\mathbf {M} ^ {\prime} \otimes_ {\mathbf {A}} \mathbf {N} ^ {\prime} \rightarrow \mathbf {M} ^ {\prime} \otimes_ {\mathbf {S} ^ {- 1} \mathbf {A}} \mathbf {N} ^ {\prime}
\]

由 A-双线性映射 \((x', y') \mapsto x' \otimes y'\)（从 \(M' \times N'\) 到 \(M' \otimes_{S^{-1}} A N'\)）导出，是双射的。

\end{enumerate}
断言 (i) 是定义 \(\mathbf{S}^{-1}\mathbf{M} =\)

\(M \otimes_{A} S^{-1} A\) 与张量积结合性的直接推论，由此在典范同构的意义下得到

\[
\begin{array}{r l} \left(\mathrm{S} ^ {- 1} \mathrm{M} \otimes_ {\mathrm{S} ^ {- 1} \mathrm{A}} \left(\mathrm{S} ^ {- 1} \mathrm{N}\right) = (\mathrm{S} ^ {- 1} \mathrm{M}) \otimes_ {\mathrm{S} ^ {- 1} \mathrm{A}} \left(\mathrm{S} ^ {- 1} \mathrm{A} \otimes \mathrm{N}\right) = (\mathrm{S} ^ {- 1} \mathrm{M}) \otimes_ {\mathrm{A}} \mathrm{N} \right. \\ = (\mathrm{S} ^ {- 1} \mathrm{A}) \otimes_ {\mathrm{A}} \mathrm{M} \otimes_ {\mathrm{A}} \mathrm{N} = \mathrm{S} ^ {- 1} (\mathrm{M} \otimes_ {\mathrm{A}} \mathrm{N}). \end{array}
\]

为证明 (ii)，我们首先注意到，在视为 A-模的 \(\mathbf{M}'\) 与 \(\mathbf{N}'\) 中，由元素 \(s \in S\) 诱导的位似映射是双射的，因而 \(\mathbf{M}' = \mathbf{S}^{-1}\mathbf{M}'\) 且 \(\mathbf{N}' = \mathbf{S}^{-1}\mathbf{N}'\)（第 2 节，注 4），类似地 \(\mathbf{S}^{-1}(\mathbf{M}' \otimes_{\mathbf{A}} \mathbf{N}') = \mathbf{M}' \otimes_{\mathbf{A}} \mathbf{N}'\)；于是 (ii) 是 (i) 的一个特殊情形。

推论. 设 M 是 A-模，a 是 A 的理想。\(\mathbf{S}^{-1}\mathbf{A}\) 的子模 \((\mathbf{S}^{-1}\mathfrak{a})(\mathbf{S}^{-1}\mathbf{M})\)、\(\mathfrak{a}(\mathbf{S}^{-1}\mathbf{M})\)、\((\mathbf{S}^{-1}\mathfrak{a})j(\mathbf{M})\)（其中 \(j: \mathbf{M} \to \mathbf{S}^{-1}\mathbf{M}\) 是典范映射）与 \(\mathbf{S}^{-1}(\mathfrak{a}\mathbf{M})\)（\(\mathbf{S}^{-1}\mathbf{M}\) 的）是相同的。特别地，若 \(\mathfrak{a}\) 与 \(\mathfrak{b}\) 是 A 的两个理想，则

\[
(S ^ {- 1} a) (S ^ {- 1} b) = a (S ^ {- 1} b) = (S ^ {- 1} a) b = S ^ {- 1} (a b).
\]

注. 设 M、N、P 为三个 A-模，\(f: M \times N \to P\) 为 A-双线性映射，而 \(S^{-1}f: (S^{-1}M) \times (S^{-1}N) \to (S^{-1}P)\) 为 \(S^{-1}A\)-双线性映射，它由 f 通过把标量环扩张到 \(S^{-1}A\) 而得到（《代数学》第 IX 章，§ 1，第 4 节，命题 1）。设 \(i: M \to S^{-1}M\)、\(j: N \to S^{-1}N\) 为典范同态；立即可得，若 Q 是由 \(f(M \times N)\) 生成的 P 的子 A-模，则 \(S^{-1}Q\) 是子 \(S^{-1}A\)-模（\(S^{-1}P\) 的），由 \((S^{-1}f)(i(M) \times j(N))\) 生成。

命题 19. 设 A 是环，S 是 A 的乘性子集。

\begin{enumerate}
\def\labelenumi{(\roman{enumi})}
\tightlist
\item
  若 M 和 N 是两个 A-模且 M 是有限生成的（相应地，有限表示的），则典范同态（第 I 章，§ 2，第 10 节，公式 (10)）

\[
\mathrm{S} ^ {- 1} \operatorname{Hom} _ {\mathrm{A}} (\mathrm{M}, \mathrm{N}) \rightarrow \operatorname{Hom} _ {\mathrm{S} ^ {- 1} \mathrm{A}} (\mathrm{S} ^ {- 1} \mathrm{M}, \mathrm{S} ^ {- 1} \mathrm{N})
\]

是单射的（相应地，双射的）。

\item
  若 \(M'\)、\(N'\) 是两个 \(S^{-1}A\)-模，则典范双射

\[
\operatorname{Hom} _ {\mathbf {S} ^ {- 1} \mathbf {A}} (\mathbf {M} ^ {\prime}, \mathbf {N} ^ {\prime}) \rightarrow \operatorname{Hom} _ {\mathbf {A}} (\mathbf {M} ^ {\prime}, \mathbf {N} ^ {\prime})
\]

是双射的。

\end{enumerate}
由于 \(S^{-1}A\) 是平坦的 A-模，(i) 是第 I 章 § 2 第 10 节命题 11 的一个特殊情形。另一方面，我们已在第 2 节命题 3 的证明过程中注意到，\(S^{-1}A\)-模的每个 A-同态必然是 \(S^{-1}A\)-同态；由此得 (ii)。

命题 20. 设 A、A' 是两个环，\(\rho: \mathbf{A} \to \mathbf{A}'\) 是同态，S 是 A 的乘性子集，\(S' = \rho(S)\)，而 \(\rho': S^{-1}A \to S'^{-1}A'\) 是对应于 \(\rho\) 的同态（第 1 节，命题 2）。

\begin{enumerate}
\def\labelenumi{(\roman{enumi})}
\tightlist
\item
  对每个 \(\mathbf{A}'\)-模 \(\mathbf{M}'\)，存在唯一的 \(\mathbf{S}^{-1}\mathbf{A}\)-同构

\[
j \colon \mathrm{S} ^ {- 1} \rho_ {*} (\mathrm{M} ^ {\prime}) \rightarrow \rho_ {*} ^ {\prime} (\mathrm{S} ^ {- 1} \mathrm{M} ^ {\prime})
\]

它满足 \(j(m' / s) = m' / \rho(s)\)（对所有 \(m' \in \mathbf{M}'\)、\(s \in \mathbf{S}\)）。

\item
  对每个 A-模 M，存在唯一的同构

\[
j ^ {\prime} \colon (\mathrm{S} ^ {- 1} \mathrm{M}) \otimes_ {\mathrm{S} ^ {- 1} \mathrm{A}} (\mathrm{S} ^ {\prime - 1} \mathrm{A} ^ {\prime}) \rightarrow \mathrm{S} ^ {\prime - 1} (\mathrm{M} \otimes_ {\mathrm{A}} \mathrm{A} ^ {\prime})
\]

（\(\mathbf{S}^{\prime -1}\mathbf{A}^{\prime}\)-模的），使得 \(j^{\prime}((m / s)\otimes (a^{\prime} / s^{\prime})) = (m\otimes a^{\prime}) / (\rho (s)s^{\prime})\)

\end{enumerate}
\begin{enumerate}
\def\labelenumi{(\roman{enumi})}
\item
  如果我们把 \(S^{\prime-1}M^{\prime}\) 借助复合同态 \(i_{M^{\prime}}^{S^{\prime}}\circ\rho\) 视为 A-模，则由 S 的元素诱导的位似映射是双射的，因而存在唯一的具有所述性质的同态 j（第 2 节，命题 3）。由于 \(\rho(S)=S^{\prime}\)，j 是满射的；此外，若 \(m^{\prime}\in M^{\prime}\)、\(s\in S\) 满足 \(m^{\prime}/\rho(s)=0\)，则存在 \(t^{\prime}\in S^{\prime}\) 使得 \(t^{\prime}m^{\prime}=0\)；由于存在 \(t\in S\) 使得 \(\rho(t)=t^{\prime}\)，于是 \(t^{\prime}.m^{\prime}=0\) 在 \(\rho_{*}(M^{\prime})\) 中成立，从而 \(m^{\prime}/s=0\) 在 \(S^{-1}\rho_{*}(M^{\prime})\) 中成立。
\item
  由于 \((\mathbf{S}^{-1}\mathbf{M})\otimes_{\mathbf{S}^{-1}\mathbf{A}}(\mathbf{S}^{\prime -1}\mathbf{A}^{\prime}) = (\mathbf{M}\otimes_{\mathbf{A}}\mathbf{S}^{-1}\mathbf{A})\otimes_{\mathbf{S}^{-1}\mathbf{A}}(\mathbf{S}^{\prime -1}\mathbf{A}^{\prime})\) 以及

\[
\mathbf {S} ^ {- 1} (\mathbf {M} \otimes_ {\mathbf {A}} \mathbf {A} ^ {\prime}) = (\mathbf {M} \otimes_ {\mathbf {A}} \mathbf {A} ^ {\prime}) \otimes_ {\mathbf {A} ^ {\prime}} (\mathbf {S} ^ {- 1} \mathbf {A} ^ {\prime}),
\]

\(j'\) 的存在性由张量积的结合性得出。

\end{enumerate}
